\section{直線配置の meridian blow-up と全 face 置換の障害}
\label{cs6:section}

境界における一つの非 toric blow-up から得られる具体的な filling を固定し、
\emph{四つの標準 pants を、互いに離した Seifert face 全体だけで貼る置換は不可能}
であることを、基本群と境界の JSJ graph によって示す。
後の定理~\ref{cs7:complexitydichotomy} と系~\ref{cs7:fourpantsnogo}では、
閉じた corner circuit に regular torus bundles を加える場合まで障害を拡張する。
open corner chains への接続などを許す一般の置換は、なお別の問題である。
したがって、以下は一般の core subdivision の否定ではなく、
実現に必要な追加構造を特定する結果である。

\subsection{一つの node を取り込む具体的なモデル}

$L_1,\ldots,L_5\subset\mathbb P^2$ を三重点を持たない五直線とし、
$p\in L_1\setminus\bigcup_{j>1}L_j$ とする。
以下では
\[
 M=\operatorname{Bl}_p\mathbb P^2,\qquad
 D'_1=H-E,\qquad D'_j=H\ (2\leq j\leq5),\qquad
 D'=\bigcup_{j=1}^5D'_j,\qquad U=M\setminus D'
\]
と書く。$H$ は直線類の引き戻し、$E$ は例外曲線である。
$K\subset U$ は十分大きな compact Stein core、$Y=\partial K$ はその正の境界とする。
$Y$ の向きは divisor neighborhood の境界 $\partial N(D')$ と逆である。

このモデルを実際の Lagrangian fibration に適用するには、取り込む node が
$L_1$ 上の非 toric blow-up に対応し、その vanishing circle が $L_1$ の meridian に
対応することを確認する必要がある。
almost toric blow-up と通常の symplectic blow-up の対応については
Chakravarthy--Groman~\cite[Theorem~1.1, Lemma~3.7]{CGblowup} を参照する。
任意の位置・任意の vanishing cycle を持つ focus--focus fiber を
Euler 数だけでこのモデルと同一視することはしない。

\begin{proposition}[一つの meridian を殺す filling]
\label{cs6:invariants}
上の $U$ に対して、次が成り立つ。
\[
 \pi_1(U)\cong\mathbb Z^3,\qquad
 H_1(U;\mathbb Z)\cong\mathbb Z^3,\qquad
 H_2(U;\mathbb Z)\cong\mathbb Z^6,\qquad
 H_k(U;\mathbb Z)=0\quad(k>2).
\]
compact core $K$ の絶対 intersection form は零であり、
$\chi(K)=4$、$\sigma(K)=0$ である。
さらに、$W_5=\mathbb P^2\setminus\bigcup_{i=1}^5L_i$ とすると、
\begin{equation}
 U\simeq W_5\cup_{\mu_1}D^2,
 \label{cs6:cellattachment}
\end{equation}
ここで $\mu_1$ は $L_1$ の正の meridian である。
\end{proposition}

\begin{proof}
$E^\circ=E\setminus D'\cong\mathbb A^1$ であり、
$U\setminus E^\circ\cong W_5$ である。
$E^\circ$ の法束は可縮な基底上の複素直線束なので自明である。
その tubular neighborhood は可縮で、例外曲線を除いた部分は $S^1$ に homotopy 同値である。
この neighborhood と $U\setminus E^\circ$ による homotopy pushout は
\eqref{cs6:cellattachment} を与える。
局所座標を $L_1=\{x=0\}$、$p=(0,0)$ とし、blow-up chart を
$x=s,y=st$ とすれば、$E^\circ$ の normal circle は $s$ の meridian であり、
blow-down 後は $\mu_1$ になる。

一般位置の五直線 complement の基本群は
\[
 \pi_1(W_5)=\langle\mu_1,\ldots,\mu_5\mid
 [\mu_i,\mu_j]=1,\ \mu_1\cdots\mu_5=1\rangle
 \cong\mathbb Z^4
\]
である。一般位置 arrangement の homotopy 型は Hattori~\cite{Hattori} による。
従って一つの primitive meridian を殺すと $\pi_1(U)=\mathbb Z^3$ になる。

$M\cong\mathbb F_1$ 上で $D'=5H-E$ は ample である。
例えば $H=E+F$ と書けば $D'=4E+5F$ であり、
Hirzebruch 曲面上の正性条件を満たす。
従って $U$ は smooth affine surface で、有限な二次元 CW complex の homotopy 型を持つ。
五つの rational curve と十個の node から
$\chi(D')=10-10=0$、一方 $\chi(M)=4$ なので $\chi(U)=4$ である。
二次元 CW complex の $H_2$ は自由群であることから $H_2(U)=\mathbb Z^6$ が従う。

$H_2(K)\to H_2(M)$ の像は全ての $D'_i$ と直交する。
$[D'_1]=H-E$ と $[D'_2]=H$ は $H_2(M)$ を張るので、この像は零である。
$K$ 内の二次元 cycle の交点数は $M$ 内でも同じであるから、絶対 intersection form は零になる。
\end{proof}

\begin{remark}[cohomology ring に残る情報]
\label{cs6:cup}
\eqref{cs6:cellattachment} の relative homology sequence では
$H_2(U,W_5)\to H_1(W_5)$ が primitive な $\mu_1$ を送る単射である。
従って $H_2(W_5)\to H_2(U)$ は同型で、$H^1(U)$ は
$H^1(W_5)$ のうち $\mu_1$ を消す rank~$3$ 部分格子である。
$W_5$ の cohomology ring が rank~$4$ exterior algebra の次数 $2$ までの切断であることから、
\[
 \bigwedge^2 H^1(U;\mathbb Z)\longrightarrow H^2(U;\mathbb Z)
\]
は primitive な rank~$3$ の像を持つ。
残りの rank~$3$ の二次 cohomology は一次 cohomology の積では生成されない。
この ring 計算だけでは homotopy 型は決まらない。次では marked cellular model から wedge decomposition 自体を証明する。
\end{remark}

\subsection{複数の meridian blow-up：homotopy 型と整数交点格子}
\label{cs7:meridian-section}

一つの node を取り込む場合の比較を、同じ成分上の複数の blow-up にも使える形にする。
この節の計算は実際の blow-up complement に対するものであり、任意の focus--focus
fiber の取り込みをこの模型で表せるという主張ではない。

$m\geq4$ 本の一般位置の直線 $L_1,\ldots,L_m\subset\PP^2$ を固定する。
$L_i$ の double points を避けた相異なる $r_i\geq0$ 点を blow up し、
得られる曲面を $M_{\boldsymbol r}$、直線の strict transforms の union を
$D_{\boldsymbol r}$、その補集合を $U_{\boldsymbol r}$ とする。
exceptional divisors は $D_{\boldsymbol r}$ に含めない。
\[
 r=\sum_i r_i,\qquad s=\#\{i:r_i>0\},\qquad
 n=m-1,\qquad k=m-1-s
\]
と書き、以下では $s<m$、すなわち少なくとも一本の直線を blow up しないと仮定する。
$T^k_{(2)}$ は標準積 CW 構造の $2$-skeleton を表し、$k\leq2$ では $T^k$ 自身とする。

\begin{proposition}[meridian attachments の homotopy 型]
\label{cs7:homotopy}
上の仮定の下で、
\begin{equation}\label{cs7:wedge}
 U_{\boldsymbol r}\simeq
 T^k_{(2)}\vee\bigvee^{\,t}S^2,
 \qquad
 t=\binom{m-1}{2}-\binom{k}{2}+r-s.
\end{equation}
従って
\[
 \pi_1(U_{\boldsymbol r})=\ZZ^k,\qquad
 b_2(U_{\boldsymbol r})=\binom{m-1}{2}+r-s.
\]
特に命題~\ref{cs6:invariants} の $U$ では
\begin{equation}\label{cs7:sextic-wedge}
 U\simeq T^3_{(2)}\vee S^2\vee S^2\vee S^2.
\end{equation}
$R=\ZZ[x_1^{\pm1},x_2^{\pm1},x_3^{\pm1}]$ とすると、その deck action を含めて
\begin{equation}\label{cs7:sextic-pi2}
 \pi_2(U)\cong R^4
\end{equation}
である。
\end{proposition}

\begin{proof}
各 exceptional divisor の $U_{\boldsymbol r}$ 内の部分は $\mathbb A^1$ であり、
これらは互いに交わらない。命題~\ref{cs6:invariants} と同じ tubular-neighborhood
pushout を各部分に適用すると、$U_{\boldsymbol r}$ は $W_m$ に $r_i$ 個ずつ
$L_i$ の meridian を attaching loop とする $2$-cells を付けた空間と homotopy 同値である。
これは相異なる点での blow-up であることを使い、infinitely near points は含めない。

blow up していない直線を infinity に選ぶと、残る $m-1$ 本の affine meridians は
Hattori の $T^{m-1}_{(2)}$ 模型の標準 generators に対応する。
ここでは単に cohomology ring を用いるのでなく、arrangement complement から
$(\CC^*)^{m-1}$ への defining-linear-functions map と整合する skeleton 模型を用いる。
この標識は Hattori~\cite[\S2, equation~(2.5), Proposition~(2.16)(2), pp.~207--213;
\S4, Lemma~(4.2), pp.~217--219]{Hattori} の構成に含まれる。
特に Proposition~(2.16)(2) は torus への inclusion との homotopy 可換性を保証し、
\S4 の local trivialization は各 labeled coordinate hyperplane を保存する。
一般位置 arrangement の skeleton 型の解説と meridian の cohomological marking については
Papadima--Suciu~\cite[\S\S1.1, 4.3; proof of Theorem~4.8(4)]{PapadimaSuciu}
も参照されたい。

この標準模型は一つの $0$-cell、$n$ 個の $1$-cells、および
$[x_i,x_j]$ を attaching words とする $\binom n2$ 個の $2$-cells を持つ。
殺す $s$ 個の generator について、それぞれ一本の circle と新しい一枚の disk を
合わせた contractible subcomplex を collapse する。
残る $k$ 個の generators の間の commutator cells は $T^k_{(2)}$ を与える。
殺した generator を含む commutator の attaching loop は null-homotopic となり、
それぞれ一つの $S^2$ を与える。同じ meridian に沿う二枚目以降の新しい cells も
同様に $S^2$ を与える。この数は $\binom n2-\binom k2+r-s$ である。
collapse した subcomplex は contractible な CW subcomplex なので、この操作は
homotopy 同値を与える。

$k=3$ のとき、$T^3$ の universal cover の cellular complex は
$x_1-1,x_2-1,x_3-1$ の Koszul complex である。
その最後の differential
\[
 d_3:R\longrightarrow R^3,\qquad
 1\longmapsto (x_3-1)e_{12}-(x_2-1)e_{13}+(x_1-1)e_{23}
\]
は単射であり、$\ker d_2=\operatorname{im}d_3\cong R$ である。
universal cover で追加の各 $S^2$ は自由 $R$ summand を一つずつ与える。
$m=5,r=s=1$ では $t=3$ なので \eqref{cs7:sextic-pi2} が従う。
\end{proof}

\begin{proposition}[整数交点格子と境界の torsion]
\label{cs7:lattice}
$K_{\boldsymbol r}\subset U_{\boldsymbol r}$ を divisor neighborhood の補集合として得られる
compact core、$Y_{\boldsymbol r}=\partial K_{\boldsymbol r}$ とする。
$H_2(K_{\boldsymbol r};\ZZ)$ の絶対 intersection lattice は
\begin{equation}\label{cs7:intersection}
 0^{\binom{m-1}{2}}\ \oplus\!
 \bigoplus_{i:r_i>0}(-A_{r_i-1}),
\end{equation}
である。$A_0=0$ とし、$A_q$ は対角 $2$、隣接成分 $-1$ の標準 root lattice とする。
特に
\[
 \sigma(K_{\boldsymbol r})=-(r-s).
\]
境界は plumbing matrix $J_m-\operatorname{diag}(r_1,\ldots,r_m)$ を持ち、
\begin{equation}\label{cs7:boundary-h1}
 H_1(Y_{\boldsymbol r};\ZZ)\cong
 \ZZ^{\binom{m-1}{2}+m-1-s}
 \oplus\bigoplus_{i:r_i>0}\ZZ/r_i\ZZ.
\end{equation}
その reduced JSJ graph は $K_m$ である。
境界包含は $H_1(Y_{\boldsymbol r})\to H_1(K_{\boldsymbol r})=\ZZ^k$ の全射を誘導し、
その kernel は
\[
 \ZZ^{\binom{m-1}{2}}\oplus\bigoplus_{i:r_i>0}\ZZ/r_i\ZZ
\]
に同型である。この抽象群の同型だけでは、境界内の具体的な peripheral subgroup は定まらない。
\end{proposition}

\begin{proof}
exceptional classes を $E_{i,a}$ と書けば
\[
 [D_i]=H-\sum_{a=1}^{r_i}E_{i,a}.
\]
pair $(M_{\boldsymbol r},U_{\boldsymbol r})$ の exact sequence と divisor に対する
duality により
\[
 \operatorname{im}\bigl(H_2(U_{\boldsymbol r})\to H_2(M_{\boldsymbol r})\bigr)
 =\ker\!\left(H_2(M_{\boldsymbol r})\xrightarrow{(\,\cdot D_i\,)_{i=1}^m}\ZZ^m\right).
\]
ここで $H_2(M_{\boldsymbol r},U_{\boldsymbol r})\cong H^2(D_{\boldsymbol r})\cong\ZZ^m$ である。
blow up していない成分との交点が零なので、kernel の $H$ coefficient は零である。
残る条件は各 $i$ について $\sum_a b_{i,a}=0$ となる。
従って kernel は
\[
 \bigoplus_{i:r_i>0}
 \left\{\sum_a b_{i,a}E_{i,a}:\sum_a b_{i,a}=0\right\},
\]
すなわち $\bigoplus_i(-A_{r_i-1})$ である。
この自由群への全射は整数加群として split し、補集合内の交点数は compactification
内の交点数と一致する。命題~\ref{cs7:homotopy} の $b_2$ から radical の rank は
$\binom{m-1}{2}$ となり、\eqref{cs7:intersection} を得る。

plumbing graph の各 edge は以前と同じ double point に対応し、各 diagonal weight は
$1-r_i$ である。よって matrix は表示されたものになる。
blow up していない一つの row を選ぶと、その relation は $\sum_i\mu_i=0$ である。
他の rows からこの row を引くと、残る relations は $r_i\mu_i=0$ となる。
未変更成分の meridian を一つ消去すれば
\[
 \operatorname{coker}Q\cong\ZZ^{m-1-s}\oplus\bigoplus_{r_i>0}\ZZ/r_i\ZZ.
\]
plumbing の exact sequence は、これに
$H_1(K_m;\ZZ)\cong\ZZ^{\binom{m-1}{2}}$ を加える。
後者は自由群なので sequence は群として split する。

各 Seifert piece の base は $(m-1)$ 穴 sphere で、その Euler 数は $3-m<0$ である。
隣接する normal-circle slopes の交点数は $1$ のままである。
従って命題~\ref{cs6:boundary} と同じ incompressibility と slope の議論により、
全ての plumbing tori が reduced JSJ system をなす。
最後の kernel の記述は、pair $(K_{\boldsymbol r},Y_{\boldsymbol r})$ の exact sequence
と $H^3(K_{\boldsymbol r})=0$ による全射、およびその自由群 $\ZZ^k$ への splitting から従う。
\end{proof}

\begin{remark}[Stein model の contact 不変量]
\label{cs7:contact}
この補集合は affine で、canonical line bundle は自明である。
実際、$N\gg0$ に対して $NH-\sum E_{i,\ell}$ は ample であり、
blow up した成分に係数 $1$、残る成分に正の整数係数を配分すれば、
この class は $D_{\boldsymbol r}$ を support とする ample divisor で表される。
また、未変更成分 $D_j=H$ を一つ選ぶと
\[
 K_{M_{\boldsymbol r}}=-\sum_{i:r_i>0}D_i+(s-3)D_j
\]
である。対応する Stein contact boundary には
\[
 c_1(\xi)=0,\qquad
 d_3(\xi)=\frac{-2\bigl(\binom{m-2}{2}+r\bigr)+3(r-s)}4
\]
が成り立つ\cite[Definition~4.2, Theorem~4.5]{Gompf}。
これも contactomorphism の十分条件ではない。
\end{remark}

\begin{corollary}[canonical pants 数に合わせる meridian 模型]
\label{cs7:canonical-budget}
$m\geq5$ とし、$P=(m-3)^2$ 個の標準四次元 pants と同じ Euler 数を持つように
$r=\binom{m-3}{2}$ 個の上記 blow-ups を行う。
さらに $\pi_1(U_{\boldsymbol r})=\ZZ^3$ を要求すると
\[
 s=m-4,\qquad
 U_{\boldsymbol r}\simeq T^3_{(2)}\vee\bigvee^{P-1}S^2,
 \qquad\pi_2(U_{\boldsymbol r})\cong\ZZ[\ZZ^3]^P,
\]
\begin{equation}\label{cs7:signature-budget}
 \sigma(K_{\boldsymbol r})=-\binom{m-4}{2}.
\end{equation}
$m\geq6$ なら、この compact core は disjoint full-face tree assembly の
$P$ 個の pants とは diffeomorphic でない。
\end{corollary}

\begin{proof}
命題~\ref{cs7:homotopy} より $m-1-s=3$ であり、同命題に $r,s$ を代入すると
wedge 数は $P-1$ となる。signature は
$r-s=\binom{m-4}{2}$ から従う。

full-face tree assembly について、face の $H_1$ からそれぞれの pants の $H_1$ への
写像は同型である。従って Mayer--Vietoris sequence を tree に沿って反復すると、
assembly の $H_2$ は各 pants に含まれる $H_2$ classes で生成される。
標準 pants の絶対交点形式は零であり、異なる pants に含まれる cycles は各 piece の
内部へ移せば互いに交わらない。よって assembly の絶対交点形式も零である。
$m\geq6$ では \eqref{cs7:signature-budget} は負なので矛盾する。
\end{proof}

例えば $m=6$、$(r_1,r_2)=(2,1)$ では $P=9$、
$U\simeq T^3_{(2)}\vee8S^2$ であるが、交点格子には $-A_1$ が現れ、
$H_1(Y)=\ZZ^{13}\oplus\ZZ/2\ZZ$ となる。
同じ homotopy 型の候補を比較しても、境界と交点格子は別に検査する必要がある。
これらの位相計算だけから Weinstein equivalence や exact symplectic surgery は従わない。


\subsection{境界 plumbing、peripheral kernel、contact plane field}

\begin{proposition}[固定すべき境界]
\label{cs6:boundary}
$\partial N(D')=-Y$ の plumbing graph は、頂点の self-intersection weight が
$(0,1,1,1,1)$ である完全 graph $K_5$ である。その plumbing matrix は
\[
 Q=J_5-\operatorname{diag}(1,0,0,0,0)
 =\begin{pmatrix}
 0&1&1&1&1\\
 1&1&1&1&1\\
 1&1&1&1&1\\
 1&1&1&1&1\\
 1&1&1&1&1
 \end{pmatrix}.
\]
この matrix は rank~$2$、Smith normal form は
$\operatorname{diag}(1,1,0,0,0)$ であり、
\begin{equation}
 H_1(Y;\mathbb Z)\cong\mathbb Z^9.
 \label{cs6:h1boundary}
\end{equation}
さらに $Y$ の reduced JSJ graph は $K_5$ そのものである。
\end{proposition}

ここでの $K_5$ は三次元境界の Seifert pieces の隣接 graph である。
Lagrangian fibration の二次元基底で core vertex の link として現れる graph とは、
異なる不変量である。

\begin{proof}
五つの strict transform は以前と同じ十個の transverse double point で交わり、
blow-up した成分だけ self-intersection が $1$ から $0$ に下がる。
これは表示された plumbing を与える。
$Q$ の第二行から第一行を引くと第一座標だけが残り、
残りは最後の四座標の和だけになるので、rank と Smith form が得られる。

$N=N(D')$ は五つの二次元 sphere を $K_5$ に従って貼った空間へ retract する。
従って $H_2(N)=\mathbb Z^5$、$H_1(N)=\mathbb Z^6$、$H_3(N)=0$ である。
Poincar\'e--Lefschetz duality と pair $(N,\partial N)$ の exact sequence により
\[
 0\longrightarrow\operatorname{coker}Q
 \longrightarrow H_1(Y;\mathbb Z)
 \longrightarrow H_1(K_5;\mathbb Z)\longrightarrow0
\]
が得られる。両端は $\mathbb Z^3$ と $\mathbb Z^6$ なので
\eqref{cs6:h1boundary} が従う。

各 Seifert piece は $S^1\times\Sigma_{0,4}$ である。
ここで $\Sigma_{0,4}$ は四穴 sphere であり、$\chi(\Sigma_{0,4})=-2$ である。
各境界 torus はこの piece に incompressible である。
plumbing gluing は一方の circle fiber と他方の circle fiber を交点数 $\pm1$ の
二つの slope に送るため、隣接する Seifert fibrations は接続しない。
また、base Euler 数が負なので Seifert fibration の曖昧さを生む
$T^2\times I$、solid torus、twisted $I$-bundle の例外は生じない。
従って十個の plumbing tori は reduced JSJ system である。
この判定は Neumann~\cite[\S5, pp.~321--322]{Neumann} の reduced graph criteria
にも直接一致する。projective line arrangement の graph-manifold の記述については
Cohen--Suciu~\cite[\S2.5]{CohenSuciu} も参照されたい。
\end{proof}

\begin{corollary}[peripheral homology の必要条件]
\label{cs6:peripheral}
境界包含 $i:Y\hookrightarrow K$ には exact sequence
\begin{equation}
 0\longrightarrow\mathbb Z^6
 \longrightarrow H_1(Y;\mathbb Z)
 \xrightarrow{i_*}\mathbb Z^3\longrightarrow0
 \label{cs6:peripheralsequence}
\end{equation}
がある。境界を固定した diffeomorphism による置換は、
境界多様体の同型だけでなく、この rank~$6$ の peripheral kernel を保存しなければならない。
\end{corollary}

\begin{proof}
$H_2(K)\to H_2(K,Y)$ は intersection map であり、命題~\ref{cs6:invariants} により零である。
$H_2(K,Y)\cong H^2(K)=\mathbb Z^6$、
$H_1(K,Y)\cong H^3(K)=0$ を pair の exact sequence に代入すればよい。
\end{proof}

plumbing fiber meridian を $\mu_i$ とすると、$H_1(Y)$ では
\[
 \mu_1=0,\qquad\mu_2+\mu_3+\mu_4+\mu_5=0.
\]
これら meridian が張る $\operatorname{coker}Q\cong\mathbb Z^3$ は
$i_*$ により $H_1(K)$ へ同型に写る。
ただし $\mu_1=0$ は homology の主張であり、boundary fundamental group で
対応する Seifert fiber が自明になるという主張ではない。

\begin{proposition}[contact plane field と log canonical 数]
\label{cs6:contact}
$U$ の canonical line bundle は代数的に自明である。
ample divisor exhaustion から得られる Stein core の contact boundary $(Y,\xi)$ は
\[
 c_1(\xi)=0,\qquad
 \theta(\xi)=-8,\qquad d_3(\xi)=-2
\]
を満たす。ここで $d_3=\theta/4$ の規約を用いる。
一方、選んだ compactification の log canonical square は
\[
 (K_M+D')^2=(2H)^2=4.
\]
\end{proposition}

\begin{proof}
$K_M=-3H+E=-D'_1-D'_2-D'_3$ であり、右辺は全て取り除いた divisor 上に支持される。
従って $K_U\cong\mathcal O_U$ である。
Stein boundary 上で $c_1(\xi)$ は $c_1(TU)$ の制限であるから零である。
Gompf の invariant の定義と不変性
\cite[Definition~4.2, Theorem~4.5]{Gompf} により
\[
 \theta(\xi)=c_1(TK)^2-2\chi(K)-3\sigma(K)=0-8-0=-8.
\]
$K_M+D'=(-3H+E)+(5H-E)=2H$ から最後の等式を得る。
\end{proof}

四直線 pants の log canonical square は $1$ なので、四つの pants はこの数とも整合する。
しかし log canonical square は固定境界の diffeomorphism 型や contact 型を決定しない。
ここでは、比較する log pair を明示しないまま log canonical square の加法性を仮定していない。
同様に $c_1(\xi)$ と $d_3(\xi)$ の一致は contactomorphism の十分条件ではない。

\subsection{全 Seifert face だけを貼る案に対する no-go}

$P$ を四直線の小さい neighborhood を $\mathbb P^2$ から取り除いた標準 compact pants とする。
その境界は四つの $S^1\times\Sigma_{0,3}$ と六つの corner tori を持ち、
boundary graph は $K_4$ である。
この記述と必要な section の選び方は
Golla--Martelli~\cite[\S4.2]{GollaMartelli} にある。

\begin{definition}[disjoint full-face assembly]
\label{cs6:assemblydef}
有限個の $P$ の互いに素なコピーを用意する。
各 boundary JSJ torus に小さい torus collar を置き、
四つの Seifert pieces を少し縮めて、その閉包が互いに離れた
$F\cong S^1\times\Sigma_{0,3}$ になるようにする。
これらの $F$ を高々一回ずつ、向きを逆にする diffeomorphism で二つずつ貼る。
残りの torus collars は外側境界に残し、角だけを丸める。
この操作以外の torus filling、handle attachment、corner 上の多重同一視、
二次元 nerve cell に対応する充填を行わないものを
\emph{disjoint full-face assembly} と呼ぶ。

各 pants を頂点、貼った face の対を辺とする graph を $\Gamma$ とする。
$\Gamma$ の loops と multiple edges は許す。
face collars が互いに離れているという条件により、全空間はこの graph 上の
通常の graph of spaces であり、隠れた triple intersection や nerve の二次元 cell を持たない。
\end{definition}

\begin{lemma}[assembly の基本群]
\label{cs6:assemblypi}
連結な disjoint full-face assembly $A$ に対して、
\[
 \pi_1(A)\cong\mathbb Z^3\rtimes_\rho F_g,
 \qquad g=b_1(\Gamma),\qquad
 \rho:F_g\longrightarrow\operatorname{GL}(3,\mathbb Z)
\]
となる。特に $\pi_1(A)\cong\mathbb Z^3$ ならば $\Gamma$ は tree である。
\end{lemma}

\begin{proof}
各標準 face の包含について
\[
 \pi_1(F)=\mathbb Z\times F_2\longrightarrow\pi_1(P)=\mathbb Z^3
\]
は可換化であり、$H_1(F)\to H_1(P)$ は整数格子の同型である。
実際、face の normal circle と、base pants の二つの boundary loops を取れば、
これらは四直線 meridian のうち三つの独立なものを、
normal meridian の整数倍を加えただけの形で与える。
従って行列式は $\pm1$ である。

van Kampen の表示では、各 edge group の commutator subgroup が両端で零になる。
face 間の diffeomorphism はこの commutator subgroup を保存するので、
同じ表示を、vertex group と edge group が全て $\mathbb Z^3$ で、
全ての edge map が同型である graph of groups として書ける。
maximal tree の辺を使って全 vertex group を一つの $\mathbb Z^3$ に同定すれば、
残りの $g$ 個の stable letters がその格子に automorphism として作用する。
これが表示された semidirect product である。

$g\geq2$ なら $F_g$ への商が非可換なので $\pi_1(A)$ は $\mathbb Z^3$ ではない。
$g=1$ でも、作用が非自明なら群が非可換で、作用が自明なら群は $\mathbb Z^4$ になる。
従って目標群が $\mathbb Z^3$ なら $g=0$ である。
\end{proof}

\begin{lemma}[tree assembly の境界 graph は planar]
\label{cs6:planar}
$\Gamma$ が tree である disjoint full-face assembly の境界は graph manifold であり、
その reduced JSJ graph は planar である。
\end{lemma}

\begin{proof}
初めに各 pants の boundary graph は planar な $K_4$ である。
異なる二つの連結成分の face を一対貼る操作は、それぞれの boundary graph から
対応する trivalent vertex を除き、三つずつの half-edge を対にする
\emph{3-edge vertex-sum} である。
三つの端点の任意の対応は、円周の順序を保つか逆にする写像で実現できる。
二つの sphere から小さい円板を除き、その境界をこの対応で貼れば再び sphere になるので、
この graph 操作は planarity を保つ。
$\Gamma$ が tree なら、葉から順にこの操作を繰り返して全境界 graph を得る。

この段階の vertex pieces は全て $S^1\times\Sigma_{0,3}$ である。
これらは irreducible で boundary-incompressible であり、
境界 tori に沿う貼合せはこの性質と貼合せ tori の incompressibility を保つ。
隣接する piece の fiber slopes が一致する torus があれば、その torus を取り除いて
二つの Seifert pieces を一つに合わせる。
base Euler 数は加法的で負のままなので、Seifert fibration の exceptional な非一意性は生じない。
base が貼合せにより非 orientable になった場合も、この負の Euler 数による一意性は同じである。
異なる fiber slopes を持つ tori だけが残った段階で reduced JSJ system となる。
graph 上ではこれらの操作は edge contraction と、同一 piece 内での loop の除去であり、
いずれも planar graph を planar graph に保つ。
\end{proof}

\begin{theorem}[全 face だけによる四 pants 置換の不可能性]
\label{cs6:nogo}
上で定義した $K$ は、disjoint full-face assembly と diffeomorphic ではない。
従って、一つの smooth boundary blow-up による
$W_5$-core と一つの node を、四つの標準 pants だけの
disjoint full-face assembly に置き換える fixed-boundary smooth surgery は存在しない。
特にこの範囲の exact symplectic surgery も存在しない。
\end{theorem}

\begin{proof}
$K\cong A$ と仮定すると、命題~\ref{cs6:invariants} と補題~\ref{cs6:assemblypi} により
gluing graph は tree である。
従って補題~\ref{cs6:planar} により $\partial A$ の reduced JSJ graph は planar である。
しかし命題~\ref{cs6:boundary} により $\partial K$ の reduced JSJ graph は非 planar な $K_5$ である。
JSJ decomposition の一意性に矛盾する。
なお、全 face の Euler 数が零なので $\chi(A)$ は pants の個数に等しく、
目標が $K$ ならその個数は自動的に四つになる。
\end{proof}

\begin{remark}[この定理が除外しない操作]
\label{cs6:nogoscope}
一般の tropical pants decomposition では、corner tori の周囲に regular torus 領域を加えて
二次元 cell を充填したり、circle collapse を行ったりする。
その場合、全空間は定義~\ref{cs6:assemblydef} の graph of spaces ではなく、
追加の二次元関係が基本群を変える。
従って補題~\ref{cs6:assemblypi} をそのまま適用できない。
Golla--Martelli の任意有限表示群の構成もこの広い設定であり、上の定理と矛盾しない。
四つの pants と Euler 数零の領域を用いる一般の置換は、依然として検討すべき問題である。
\end{remark}

\subsection{relative mapping class と monodromy の検査項目}

Golla--Martelli~\cite[Proposition~4.3]{GollaMartelli} の記述は、
corner を含めた貼合せを次に試す際の具体的な必要条件を与える。
section と fiber の基底を固定した face gluing の三つの境界行列は
\[
 L_j=\begin{pmatrix}\epsilon&k_j\\0&\epsilon'\end{pmatrix},
 \qquad \epsilon,\epsilon'\in\{1,-1\},\qquad k_1+k_2+k_3=0
\]
を満たさなければならない。
さらに、各 disk face の周囲には、vertex での基底変更行列 $J_j$ を含めて
\[
 J_nL_n\cdots J_1L_1=I
\]
が必要である。
これらは「三つの torus に別々に都合のよい行列を割り当てる」操作が
一つの $S^1\times\Sigma_{0,3}$ の diffeomorphism から来るための制約である。
face が disk でない場合は、境界 monodromy が内部へ延長することと、
その延長自体の選択がさらに必要になる
\cite[Remark~4.9]{GollaMartelli}。
これらの topological 条件に加え、Lagrangian fibration の surgery では
実際の symplectic collar、period lattice、action primitive の相対的一致が必要である。

五直線モデルからは、補助的な braid monodromy の検査も得られる。
以下は deleted divisor の monodromy であり、focus--focus monodromy と混同しない。

\begin{proposition}[ruling に付随する boundary braid]
\label{cs6:braid}
$p$ をさらに一般に選ぶと、projection from $p$ は
\[
 U\cong\bigl(\mathbb A^1_t\times\mathbb P^1_y\bigr)
 \setminus\bigcup_{i=2}^5\{y=a_it+b_i\}
\]
と書ける。$a_i$ は相異なり、六つの pairwise collision parameters は相異なる。
affine fiber chart $\mathbb A^1_y$ における四つの punctures の boundary braid は
\[
 \beta_1\cdots\beta_6=\Delta_4^2=(\sigma_1\sigma_2\sigma_3)^4\in B_4,
\]
各 $\beta_j$ は二点の正の full twist の共役である。
factorization は distinguished paths の変更で Hurwitz move と simultaneous conjugation を受ける。
\end{proposition}

\begin{proof}
$p=[0:0:1]$、$L_1=\{X=0\}$ と取る。
projection parameter を $t=Y/X$、fiber coordinate を $y=Z/X$ とすれば、
残り四直線は $y=a_it+b_i$ であり、例外曲線は $y=\infty$ の section になる。
三重点がないので $a_i$ は相異なる。
$p$ を有限個の追加の位置から避ければ、異なる pair の collision values も相異なる。

一つの collision value の周囲では、二つの punctures の差が
非零定数倍の $t-t_0$ である。正の一周でその差は一回回転するので、
local braid は正の full twist になる。
十分大きい $|t|=R$ 上では $b_i$ を零へ縮めても punctures は衝突しない。
従って boundary braid は四点 $a_it$ の同時の正の一回転、すなわち $\Delta_4^2$ である。
\end{proof}

この ruling は非 proper な submersion であり、collision は取り除いた divisor 上で起こる。
従ってここで「六つの critical values」は interior Lefschetz singularities の数を意味しない。
ruling と affine fiber chart を保持する置換なら上の factorization を合わせる必要があるが、
任意の smooth surgery に ruling の保存を要求しているわけではない。

\subsection{very affine realization に対する障害}

\begin{proposition}[例外 affine line と unit map]
\label{cs6:veryaffine}
$U$ は very affine ではない。
より強く、任意の algebraic morphism $U\to(\mathbb C^*)^N$ は
$E^\circ\cong\mathbb A^1$ 上で定数になる。
unit group は
\[
 \mathcal O(U)^*/\mathbb C^*\cong
 \{(c_1,\ldots,c_5)\in\mathbb Z^5:
 c_1=0,\ c_2+c_3+c_4+c_5=0\}\cong\mathbb Z^3
\]
である。
\end{proposition}

\begin{proof}
invertible regular function を $\mathbb A^1$ に制限すると、
零点を持たない多項式なので定数である。
torus への morphism の全座標がこの性質を持つため、$E^\circ$ を分離できない。
unit group の式は divisor exact sequence から得られる。
すなわち $M$ 上の rational function が $U$ 上で unit であることは、
その divisor が $D'_1,\ldots,D'_5$ 上に支持されることと同値であり、
\[
 \sum_i c_i[D'_i]
 =\Bigl(\sum_i c_i\Bigr)H-c_1E=0
\]
が必要十分である。
\end{proof}

この計算により、固定した複素曲面 $U$ を一つの algebraic torus 内の smooth hypersurface
としてそのまま Newton polytope で subdivision する方法は使えない。
これは別の複素構造への変形、symplectic decomposition、部分領域の very affine chart、
または non-toric な追加 chart を用いる構成の否定ではない。
$U$ は四直線 complement に対し、五本目の直線との交わりを divisor、$p$ を center とする
affine modification として記述できる。
この操作の代数的枠組みは Kaliman--Zaidenberg~\cite[\S\S1--3]{KalimanZaidenberg} にある。

